\chapter{A reproducible computational workshop}
\label{ch:workshop}

\begin{learningobjectives}
After this chapter, you should be able to:
\begin{itemize}
  \item distinguish a terminal, shell, Python interpreter, virtual environment,
  VS Code, Jupyter, and a kernel;
  \item explain what the course setup command changes and what it does not;
  \item verify the interpreter used by a notebook; and
  \item diagnose environment failures from evidence before reinstalling tools.
\end{itemize}
\end{learningobjectives}

\begin{chapterconnection}
The overview named the full evidence-to-product system. This chapter begins at
the smallest operational boundary: one student, one repository, and one Python
process. Before reasoning about data, we must know which program is executing
the code and how another machine can reproduce that choice.
\end{chapterconnection}

\section{A project has a place and a runtime}

A reproducible result depends on both source files and the software that
interprets them. Two students may open identical notebooks yet run different
Python versions or package sets. One notebook may use the project environment;
another may silently use a system interpreter.

\begin{definitionbox}{repository}
A repository is the project directory whose version-controlled files describe
the work. It includes source, notebooks, tests, configuration, and documentation.
The repository is not the same object as a Python installation or virtual
environment.
\end{definitionbox}

A \term{path} identifies a location in a filesystem. Absolute paths begin from
a filesystem root and often contain machine-specific details. Relative paths
begin from an agreed working directory. Course code should discover the project
root from stable markers such as \code{pyproject.toml}, accept paths from
callers, or use package resources. Hard-coding a student's home directory is
not portable.

\section{Terminal, shell, and Python}

These names refer to different layers:

\begin{description}
  \item[Terminal.] A text-oriented application that displays a session and
  sends keystrokes to a shell or another command-line program.
  \item[Shell.] A program such as PowerShell, Bash, or Zsh that parses commands,
  locates programs, expands variables, and connects input and output.
  \item[Python interpreter.] The executable process that parses and runs Python
  code.
  \item[VS Code.] An editor and development environment that can host a terminal,
  edit notebooks, and connect extensions to Python tools.
\end{description}

Typing \code{python analysis.py} asks the shell to locate a program named
\code{python} and pass it the path \code{analysis.py}. The shell is not Python,
and Python does not interpret shell syntax.

\begin{warningbox}
Copying an activation command for the wrong shell is a common failure. A Bash
command such as \code{source .venv/bin/activate} is not PowerShell syntax. The
course uses \code{uv run} so the essential workflow does not depend on students
memorizing shell-specific activation.
\end{warningbox}

\section{Virtual environments solve project isolation}

A \term{virtual environment} is a project-specific Python environment with its
own interpreter-facing package installation area. It answers:

\begin{quote}
Which Python and package versions should this project use without treating the
whole computer as one shared project?
\end{quote}

It does not emulate a new operating system, guarantee identical hardware, or
make arbitrary code safe. It is dependency isolation, not a security sandbox.

The repository declares intent in \code{pyproject.toml}. The lockfile records a
resolved dependency graph. The local \code{.venv} realizes that graph on one
machine. These are related but distinct artifacts.

\begin{center}
\begin{tikzpicture}[node distance=7mm]
  \node[wideflowbox] (project) {\code{pyproject.toml}\\declared requirements};
  \node[wideflowbox,right=of project] (lock) {\code{uv.lock}\\resolved versions};
  \node[wideflowbox,right=of lock] (venv) {\code{.venv}\\local installation};
  \node[wideflowbox,below=of lock] (process) {Python process\\imports and runs code};
  \draw[flowarrow] (project) -- (lock);
  \draw[flowarrow] (lock) -- (venv);
  \draw[flowarrow] (venv) -- (process);
\end{tikzpicture}
\end{center}

\section{Jupyter and the kernel}

A notebook file stores cells, metadata, and possibly outputs. It does not run
itself. A \term{kernel} is a long-running language process that receives code
from a notebook interface, retains state, and returns results.

VS Code's notebook editor is the visible interface. The Jupyter extension sends
cells to the selected kernel. The kernel runs a specific Python executable.
Selecting the wrong kernel therefore produces a characteristic puzzle: a
package is installed in one environment but cannot be imported by the Python
process actually serving the notebook.

\begin{lstlisting}[caption={Inspect the Python process serving a notebook.}]
import sys
from pathlib import Path

print("Python version:", sys.version.split()[0])
print("Python executable:", Path(sys.executable))
\end{lstlisting}

The important evidence is \code{sys.executable}. An activated terminal can use
one Python while an already-running kernel continues using another. Restarting
the kernel replaces the process and clears its in-memory state.

\section{The one-command course setup}

From the repository root in VS Code's integrated terminal, run:

\begin{lstlisting}[language=bash,numbers=none]
uv run python scripts/setup_course.py
\end{lstlisting}

This workflow synchronizes the project environment, verifies that the
\code{rice\_dsm} package imports from this repository, and registers a named
\textbf{Rice DSM} kernel. It does not launch JupyterLab, edit a student's
notebook, upload code, or install packages globally.

After setup, open a notebook in VS Code and select the Rice DSM kernel. Then
verify:

\begin{lstlisting}
import rice_dsm
import sys

print(rice_dsm.__version__)
print(rice_dsm.__file__)
print(sys.executable)
\end{lstlisting}

The package location should point into this project, and the executable should
belong to the project environment.

\section{Diagnose before reinstalling}

When an import fails, collect evidence in this order:

\begin{enumerate}
  \item Confirm the repository root: does \code{pyproject.toml} exist here?
  \item Inspect \code{sys.executable} in the failing notebook or process.
  \item Compare that path with the project's \code{.venv}.
  \item Verify the selected VS Code kernel name.
  \item Run the course setup command and read its complete output.
  \item Restart the kernel and run the notebook from the first cell.
\end{enumerate}

Repeated installation can make the system harder to understand by creating
additional environments. Diagnosis should identify the process and boundary
before changing state.

\begin{checkpointbox}
In your own words, distinguish Python, a virtual environment, Jupyter, and a
kernel. Then explain why installing a package in one environment does not make
it available to a kernel using another interpreter.
\end{checkpointbox}

\section{Worked diagnosis: an import that succeeds in the terminal only}

Imagine that Maya runs the following command in VS Code's terminal:

\begin{lstlisting}[language=bash,numbers=none]
uv run python -c "import rice_dsm; print(rice_dsm.__file__)"
\end{lstlisting}

It succeeds. In a notebook, however, \code{import rice\_dsm} raises
\code{ModuleNotFoundError}. It is tempting to conclude that installation is
random or broken. The evidence supports a more specific hypothesis: the two
commands reached different Python processes.

\begin{workedexample}{follow the processes, not the icons}
\textbf{Step 1: identify the successful process.} The shell sees \code{uv run},
which resolves the project and starts Python inside the synchronized project
environment.

\textbf{Step 2: identify the failing process.} The notebook editor sends its
cell to whichever kernel is selected in the upper-right kernel control. That
kernel may have started hours earlier from a system Python installation.

\textbf{Step 3: measure.} In the notebook, run only:

\begin{lstlisting}
import sys
print(sys.executable)
\end{lstlisting}

If the printed path does not belong to the repository's \code{.venv}, the
failure is explained. No package reinstall is needed.

\textbf{Step 4: correct the boundary.} Select the Rice DSM kernel, restart it,
and run from the first cell. A restart matters because changing a menu setting
does not transform an already-running Python process into another interpreter.

\textbf{Step 5: verify the correction.} Inspect \code{sys.executable},
\code{rice\_dsm.\_\_file\_\_}, and the package version. Verification turns a
plausible fix into evidence.
\end{workedexample}

This pattern generalizes. When two interfaces disagree, identify the actual
process, configuration, working directory, and input used by each. Avoid
reasoning from the fact that both windows ``look like Python.''

\begin{misconceptionbox}
\textbf{Misconception:} activating an environment permanently changes Python on
the computer.

\textbf{Correction:} activation usually changes shell variables so commands in
that shell locate one interpreter first. It does not rewrite other terminals,
running kernels, editor processes, or the operating system's Python. \code{uv
run} selects the project environment for the command it launches.
\end{misconceptionbox}

\section{State explains surprising notebook behavior}

A notebook kernel retains definitions and objects between cell executions. If a
student runs cell 12 before cell 4, cell 12 may succeed because a name remains
from yesterday, or fail because a required definition has not yet executed.
The notebook file's visual order and the kernel's historical execution order
are not necessarily the same.

This is why ``Restart Kernel and Run All'' is an engineering check. It asks
whether the document describes a complete computation from a fresh process.
Clearing visible outputs without restarting does not clear memory. Restarting
without running from the beginning removes hidden state but does not recreate
it.

\conceptcheck{A notebook displays output below every cell, but a fresh Run All
fails at cell 7. Which evidence should be trusted when deciding whether the
notebook is reproducible? Explain why.}

\section{Exercises}

\subsection*{Foundational}

\begin{enumerate}
  \item Draw separate boxes for VS Code, its integrated terminal, the shell, a
  Python interpreter, a notebook file, and a kernel. Add arrows showing who
  launches or communicates with whom.
  \item Classify each item as source declaration, resolved plan, or local
  realization: \code{pyproject.toml}, \code{uv.lock}, and \code{.venv}.
  \item Explain why an absolute path copied from one student's computer is
  unlikely to work for another student.
\end{enumerate}

\subsection*{Applied}

\begin{enumerate}
  \item Record \code{Path.cwd()}, \code{sys.executable}, Python's version, and
  \code{rice\_dsm.\_\_file\_\_} from a working course notebook. Annotate what
  each observation establishes and what it does not establish.
  \item Deliberately select a different available kernel, observe the evidence,
  and return to Rice DSM. Do not install anything. Write a three-sentence
  diagnosis based only on paths and process identity.
\end{enumerate}

\subsection*{Design}

Write a troubleshooting guide for a classmate who reports only ``Python is not
working.'' Your guide must begin with read-only observations, branch based on
evidence, and avoid global installation commands.

\begin{chapterreview}
\textbf{Key ideas.} A repository, environment, interpreter, editor, notebook,
and kernel are distinct objects. Reproducibility requires knowing which process
ran the code. Diagnose identity, path, and state before changing installation.

\textbf{Vocabulary.} repository; path; terminal; shell; interpreter; virtual
environment; dependency; lockfile; notebook; kernel; working directory.

\textbf{Explain aloud.} Why can \code{pip install} report success while the
next notebook import fails?
\end{chapterreview}

\begin{notebooklink}
Work through Lecture 1 notebook 00 and the computing-foundations guides. The
repository tests also exercise setup, imports, the named kernel, and every
notebook from top to bottom.
\end{notebooklink}

\section{Takeaway}

Reproducibility starts by naming the project, interpreter, dependencies, and
process that produced a result. Tools become easier to debug when each is given
one precise responsibility.
